\chapter{Proportionnalité, ratios et pourcentages}
\label{exercices:07}
\section*{Exercice 7.1 -- Reconnaître le modèle}
Les tableaux $(x;y)=(2;5),(4;10),(7;17{,}5)$ et $(2;7),(4;11),(7;17)$ sont-ils proportionnels ? Donner une formule simple pour chaque tableau.

\section*{Exercice 7.2 -- Partager et mélanger}
Partager $360$ euros selon le ratio $2:3:4$. Dans un mélange sirop:eau de ratio $1:5$, quelle quantité de sirop faut-il pour $1{,}8\unite L$ de mélange ?

\section*{Exercice 7.3 -- Évolutions successives}
Un prix de $80$ euros augmente de $15\,\%$, puis baisse de $20\,\%$. Calculer le prix final et le taux global. Quel taux de hausse permettrait ensuite de revenir à $80$ euros ?

\section*{Exercice 7.4 -- Échelle et précision}
Sur une carte au $1:50000$, une distance mesure $7{,}2\unite{cm}$. Quelle est la distance réelle en kilomètres ? Une précision de lecture de $1\unite{mm}$ correspond à quelle distance réelle ?
